% =====================================================================
%  PARTIE II : exemple multi-langages (C, SQL) + figure
% =====================================================================
\chapter{Tri d'un tableau et interrogation d'une base de données}

\section{Énoncé de l'exercice}

\begin{remarque}[Consigne]
Écrire un programme C qui trie un tableau d'entiers par la méthode du tri à
bulles, puis une requête SQL qui sélectionne les employés dont le salaire
dépasse un seuil donné.
\end{remarque}

\section{Implémentation en langage C}

Le tri à bulles parcourt le tableau et échange deux éléments consécutifs
lorsqu'ils sont mal ordonnés. Nous utilisons la fonction \codeinline{swap}
pour l'échange.

\begin{code}[language=C, numbers=left]{tri.c}
#include <stdio.h>
#include <stddef.h>

static void swap(int *a, int *b)
{
    int tmp = *a;
    *a = *b;
    *b = tmp;
}

void tri_bulles(int *tab, size_t n)
{
    for (size_t i = 0; i < n - 1; ++i) {
        for (size_t j = 0; j < n - 1 - i; ++j) {
            if (tab[j] > tab[j + 1]) {
                swap(&tab[j], &tab[j + 1]);
            }
        }
    }
}

int main(void)
{
    int tab[] = {5, 2, 9, 1, 7};
    size_t n = sizeof(tab) / sizeof(tab[0]);

    tri_bulles(tab, n);

    for (size_t i = 0; i < n; ++i) {
        printf("%d ", tab[i]);
    }
    printf("\n");
    return 0;
}
\end{code}

Compilation et exécution~:

\begin{terminal}
$ gcc -Wall -o tri tri.c
$ ./tri
1 2 5 7 9
\end{terminal}

\section{Requête SQL associée}

\begin{code}[language=SQL]{employes.sql}
SELECT nom, prenom, salaire
FROM employes
WHERE salaire > 8000
ORDER BY salaire DESC;
\end{code}

Le résultat de la requête peut être présenté sous forme de capture d'écran
(Figure~\ref{fig:capture}).

\begin{figure}[h]
  \centering
  \includegraphics[width=0.7\linewidth]{images/capture_exemple.png}
  \caption{Exemple de résultat obtenu après exécution.}
  \label{fig:capture}
\end{figure}

\section{Analyse et complexité}

Le tri à bulles a une complexité temporelle en $O(n^2)$ dans le pire cas, ce
qui le rend peu efficace pour de grands tableaux. La requête SQL, quant à
elle, s'appuie sur l'opérateur de comparaison \codeinline{salaire > 8000} et
sur la clause \codeinline{ORDER BY} pour trier le résultat.
